%% fig-execution.tex -- Figure (single column): the execution and control model.
%%
%% The claim this figure has to support: arbitration is durable where it is
%% claim-based, and per-process where it is in-memory. Those two are drawn
%% differently on purpose, because conflating them is an operational hazard.
%%
%% Every row is placed from the previous row's anchor rather than at a fixed y, so
%% the layout survives a change of wording or of type scale.

\begin{figure}[tbp]
\centering
\begin{cbscaled}[0.68]
\begin{tikzpicture}[x=1cm,y=1cm]

  \def\W{8.02}   % content width; the side lanes add ~0.3cm
  \tikzset{erow/.style = {anchor=north west, font=\figB, align=left}}

  % ---- request path --------------------------------------------------------
  \node[surface, anchor=north west, text width=\W cm-8pt, minimum height=0.42cm]
    (req) at (0,0)
    {\textbf{Request path}\ \ {\figB async routes, synchronous ORM}};

  \node[control, erow, text width=3.76cm] (inl)
    at ([yshift=-4.4mm]req.south west)
    {\textbf{Inline}\\bounded validation and\\most durable mutations};
  \node[async, erow, text width=3.76cm] (enq) at ([xshift=4.26cm]inl.north west)
    {\textbf{Enqueued}\\long-running or\\explicitly queued work};
  \draw[flow] (req.south) -- (req.south |- inl.north);

  % ---- durable queue -------------------------------------------------------
  \node[db, erow, text width=\W cm-8pt] (q) at ([yshift=-4.4mm]inl.south west)
    {\textbf{Durable queue}\ \ status and claim columns on relational rows; leases
     only where the loop defines one --- no broker in the arbitration path};
  \draw[flow, draw=cbAsync] (enq.south) -- (enq.south |- q.north);

  % ---- the loops, split by arbitration discipline --------------------------
  \node[bandlab, anchor=north west] (l1)
    at ([yshift=-3.8mm, xshift=1pt]q.south west)
    {CLAIM-ARBITRATED \normalfont\itshape --- one owner at a time};

  \node[async, erow, text width=\W cm-8pt] (w1)
    at ([yshift=-1.4mm, xshift=-1pt]l1.south west)
    {\textbf{Refinement}\hfill\textsc{atomic claim}\\
     triage, evidence, diagnosis, candidate, evaluation};
  \node[async, erow, text width=\W cm-8pt] (w2) at ([yshift=-1.1mm]w1.south west)
    {\textbf{CalibrationDrain}\hfill\textsc{atomic claim}\\
     queued calibration jobs};
  \node[async, erow, text width=\W cm-8pt] (w3) at ([yshift=-1.1mm]w2.south west)
    {\textbf{WorkflowRun}\hfill\textsc{atomic claim}\\
     queued runs, checkpoints, runtime approvals};
  \node[async, erow, text width=\W cm-8pt] (w4) at ([yshift=-1.1mm]w3.south west)
    {\textbf{KnowledgeBase}\hfill\textsc{atomic claim}\\
     ingest, chunk, embed, graph extraction};
  \node[async, erow, text width=\W cm-8pt] (w5) at ([yshift=-1.1mm]w4.south west)
    {\textbf{WebhookDispatcher}\hfill\textsc{atomic claim}\\
     delivery, settlement, dead letters};

  \node[bandlab, anchor=north west] (l2)
    at ([yshift=-3.8mm, xshift=1pt]w5.south west)
    {IDEMPOTENT OR UNLEASED \normalfont\itshape --- correct by another argument};

  \node[muted, erow, text width=\W cm-8pt] (u1)
    at ([yshift=-1.4mm, xshift=-1pt]l2.south west)
    {\textbf{WorkflowScheduler}\hfill\textsc{unique partial index}\\
     cron triggers per deployment alias};
  \node[muted, erow, text width=\W cm-8pt] (u2) at ([yshift=-1.1mm]u1.south west)
    {\textbf{Janitor}\hfill\textsc{idempotent sweep}\\
     reaps jobs past their heartbeat timeout};
  \node[muted, erow, text width=\W cm-8pt] (u3) at ([yshift=-1.1mm]u2.south west)
    {\textbf{AriaPlan}\hfill\textsc{no claim lease}\\
     resumes plans parked on async jobs};

  % the queue feeds the loop groups; status transitions flow back
  \draw[flow, draw=cbAsync] ([xshift=-3pt]q.south west)
    -- ([xshift=-3pt]w1.north west) -- ([xshift=2pt]w1.north west);
  \draw[flow, draw=cbAsync] ([xshift=3pt]u1.north east)
    -- ([xshift=3pt]q.south east) -- ([xshift=-2pt]q.south east);

  % ---- the per-process caveat ---------------------------------------------
  \node[govern, erow, text width=\W cm-8pt] at ([yshift=-4.8mm]u3.south west)
    {\textbf{Per-process, \emph{not} durable}\\[0.2ex]
     The failed-login throttle and the API rate limiter are in-memory token
     buckets behind a thread lock. Every process runs its own full set of loops,
     so with \code{mlflow server}'s four default workers these budgets exist four
     times over and the effective ceiling scales with the worker count. Size them
     accordingly, or pin the process to a single worker.};

\end{tikzpicture}
\end{cbscaled}
\caption{The execution model. Bounded validation and most durable mutations run
inline in request handlers; explicitly queued or long-running work goes to up to
\statLoops{} in-process loops, with no separate worker tier. The split between the
two loop groups is the operationally important part. Claim-based consumers compete
for rows atomically (\algref{alg:claim}), so no two workers ever hold the same row
--- but their \emph{recovery} policies differ, and \tabref{tab:loops} gives the
resulting delivery semantics loop by loop. The scheduler is idempotent by a minute-bucketed key backed by a unique
partial index, which makes duplicate fires impossible rather than merely
unlikely; the janitor is an idempotent sweep; \aria's plan worker polls without a
claim lease. The bottom box is the one place where the model is genuinely
per-process, and it is stated here rather than left to be discovered in
production.}
\label{fig:execution}
\end{figure}
